\section{Glossary}

\begin{description}[leftmargin=0.3in,style=nextline]
  \item[Finite field.]
        A finite set with addition, subtraction, multiplication, and division by
        nonzero elements.
  \item[MLE.]
        Multilinear extension: the unique multilinear polynomial extending a
        Boolean tensor to all field points.
  \item[RAA/PRAA.]
        The packed repeat-accumulate-accumulate code used by Blaze to encode
        each message row.
  \item[Merkle commitment.]
        A hash-tree commitment that lets the prover open selected flat-vector
        entries with short authentication paths.
  \item[IOP.]
        Interactive oracle proof: an interactive proof template where the prover
        fixes strings behind message-oracle handles and the
        verifier later opens selected coordinates through the handle
        interface.  Its soundness promise is exact membership in a relation.
  \item[IOPP.]
        Interactive oracle proof of proximity: a protocol realization of
        \(\mathcal F_{\mathrm{Prox}}^{L,\delta}\) for some language
        \(L\subseteq\Sigma^D\).  The verifier receives a handle to a large
        implicit input, opens selected coordinates through that handle, and
        rejects inputs farther than \(\delta\) from \(L\).  Close-but-invalid
        inputs inside the radius are in the proximity gap.
  \item[Holographic IOP.]
        An IOP for an indexed relation \(R(i,x,w)\) where the verifier receives
        a handle to preprocessed public relation data instead of reading all of
        \(i\).  Its soundness promise is exact membership for the selected
        relation, not proximity of the index handle.
  \item[MLIOP.]
        Multilinear polynomial interactive oracle proof: a polynomial-oracle
        proof template where the prover registers Boolean tensors with an ideal
        evaluation functionality that answers MLE evaluations.
  \item[Sumcheck.]
        A protocol that reduces a large hypercube sum claim to one random
        terminal check.
  \item[BaseFold.]
        A fold-and-query proof architecture whose concrete protocol realizes
        proximity for committed codeword oracles.  When superimposed with
        sumcheck, it helps authenticate terminal \(\mathsf{Eval}\) claims using
        coordinate \(\mathsf{Open}\) handles.
  \item[Holographic row.]
        A row used by the backend but authenticated by an external handle;
        here, $c_\star$ is authenticated by outer Blaze columns.
  \item[Residual.]
        A row whose entries are zero when a local relation is satisfied.
  \item[Terminal vector.]
        The five final folded values of
        $(c_\star,u_2,u_3,u_4,R_\eta)$ after all shared folds.
\end{description}

\section{References}

\begin{itemize}
  \item[\textbf{[Blaze24]}] Brehm, Chen, Fisch, Resch, Rothblum, and Zeilberger,
        ``Blaze: Fast SNARKs from Interleaved RAA Codes,''
        IACR ePrint 2024/1609.
        \url{https://eprint.iacr.org/2024/1609}
  \item[\textbf{[BaseFold23]}] Zeilberger, Chen, and Fisch,
        ``BaseFold: Efficient Field-Agnostic Polynomial Commitment Schemes from
        Foldable Codes,'' IACR ePrint 2023/1705.
        \url{https://eprint.iacr.org/2023/1705}
  \item[\textbf{[BCS16]}] Ben-Sasson, Chiesa, and Spooner,
        ``Interactive Oracle Proofs,''
        Theory of Cryptography Conference 2016.
  \item[\textbf{[RVW13]}] Rothblum, Vadhan, and Wigderson,
        ``Interactive Proofs of Proximity: Delegating Computation in
        Sublinear Time,'' STOC 2013.
  \item[\textbf{[COS20]}] Chiesa, Ojha, and Spooner,
        ``Fractal: Post-Quantum and Transparent Recursive Proofs from
        Holography,'' EUROCRYPT 2020.
  \item[\textbf{[CBBZ23]}] Chen, Bunz, Boneh, and Zhang,
        ``HyperPlonk: Plonk with Linear-Time Prover and High-Degree Custom
        Gates,'' IACR ePrint 2022/1355.
  \item[\textbf{[LFKN92]}] Lund, Fortnow, Karloff, and Nisan,
        ``Algebraic Methods for Interactive Proof Systems,'' JACM 1992.
  \item[\textbf{[CCH19]}] Canetti, Holmgren, Jain, Richelson, and Vaikuntanathan,
        ``Fiat-Shamir: From Practice to Theory,'' STOC 2019.
  \item[\textbf{[CMS19]}] Chiesa, Manohar, and Spooner,
        ``Succinct Arguments in the Quantum Random Oracle Model,''
        TCC 2019.
  \item[\textbf{[RZR21]}] Ron-Zewi and Rothblum,
        ``Local Proofs Approaching the Witness Length,'' January 2021.
        This is the local code-switching source document in the docs
        directory.
  \item[\textbf{[Walkthrough]}] Local source document:
        the full Blaze protocol walkthrough in the docs directory.
\end{itemize}
